\documentclass[11pt]{article}
\usepackage[a4paper,margin=1in]{geometry}
\usepackage[T1]{fontenc}
\usepackage{helvet}
\renewcommand{\familydefault}{\sfdefault}
\usepackage{xcolor}
\usepackage{booktabs}
\usepackage{array}
\usepackage{colortbl}
\usepackage{float}
\usepackage{amsmath,amssymb}
\usepackage[font=footnotesize,labelformat=empty]{caption}
\usepackage{titlesec}
\usepackage{fancyhdr}
\usepackage{enumitem}
\usepackage{lastpage}
\usepackage[hidelinks]{hyperref}

% ---- Bayesian Capital brand palette ----
\definecolor{bcnavy}{RGB}{11,31,58}
\definecolor{bcblue}{RGB}{32,74,135}
\definecolor{bcgold}{RGB}{176,141,87}
\definecolor{bcgrey}{RGB}{90,98,112}
\definecolor{bcrule}{RGB}{210,214,220}
\definecolor{bckeep}{RGB}{232,237,244}
\definecolor{bcact}{RGB}{247,241,228}
\definecolor{bcbad}{RGB}{247,230,230}

\titleformat{\section}{\color{bcnavy}\Large\bfseries}{\thesection}{0.6em}{}
\titleformat{\subsection}{\color{bcblue}\large\bfseries}{\thesubsection}{0.6em}{}
\titlespacing*{\section}{0pt}{1.4em}{0.5em}

\pagestyle{fancy}\fancyhf{}
\renewcommand{\headrulewidth}{0pt}
\renewcommand{\footrulewidth}{0.4pt}
\renewcommand{\footrule}{\hbox to\headwidth{\color{bcrule}\leaders\hrule height \footrulewidth\hfill}}
\fancyfoot[L]{\footnotesize\color{bcgrey}Bayesian Capital}
\fancyfoot[C]{\footnotesize\color{bcgrey}Confidential --- Internal Research}
\fancyfoot[R]{\footnotesize\color{bcgrey}Page \thepage\ of \pageref{LastPage}}

\newcommand{\kpi}[2]{\textbf{#1}\,{\footnotesize\color{bcgrey}#2}}
\setlength{\parindent}{0pt}\setlength{\parskip}{0.55em}

\begin{document}

% ================= title =================
\begin{titlepage}
\vspace*{1.1cm}
{\color{bcgold}\rule{\textwidth}{2pt}}\\[0.4cm]
{\color{bcnavy}\Huge\bfseries Execution-Verified Performance}\\[0.3cm]
{\color{bcnavy}\Huge\bfseries of the Hybrid Harvester}\\[0.35cm]
{\color{bcblue}\LARGE White Paper}\\[0.5cm]
{\color{bcgold}\rule{\textwidth}{1pt}}\\[0.6cm]
{\large\color{bcgrey} Intraday validation of same-day round trips, restated returns,}\\[0.15cm]
{\large\color{bcgrey} capital allocation, and a register of rejected hypotheses}\\[1.1cm]
{\color{bcnavy}\large\bfseries Bayesian Capital}\\[0.15cm]
{\color{bcgrey} Quantitative Research}\\[0.15cm]
{\color{bcgrey} 3 August 2026 --- Revision 2}\\[1.6cm]

\fbox{\parbox{0.92\textwidth}{\color{bcnavy}\textbf{Abstract.}\\ \color{black}
The daily-bar backtest of the hybrid Bayesian~/~Ornstein--Uhlenbeck harvester cannot determine
whether an intraday low preceded the corresponding high, and therefore resolves every same-day
round trip in the strategy's favour. Using minute and five-minute data for all ten names we
find that only $57\%$ of such trades were actually achievable. Restating the book on verified
fills reduces the mean annualised return from $\sim\!220\%$ to $\sim\!45\%$. The edge is real
and substantial; its magnitude was overstated roughly fivefold. Nine changes were tested out of
sample and seven failed. The two that survived are both allocation decisions resting on
measured, persistent properties of the data: equal weighting across names, and a tilt of
capital toward the Bayes sleeve. Every attempt to fit better parameter values failed.}}
\end{titlepage}

% ================= 1 =================
\section{Introduction and scope}

The harvester allocates capital across ten high-beta US equities through two decorrelated
sleeves: a Bayesian sleeve driven by a Kalman local-linear-trend filter, and a mean-reverting
sleeve driven by an Ornstein--Uhlenbeck (OU) process. Both place resting limit bids below a
model-derived fair value and exit at a fixed take-profit premium or a calendar stop.

All performance figures to date derive from a daily-bar backtest. This paper asks whether those
figures are attainable, using intraday data to test the one assumption daily bars cannot
verify. It then restates the book's expected performance, re-examines capital allocation on the
corrected basis, and records the hypotheses tested and rejected.

All computations use a Python replica of the workbook validated to reproduce every cached
\texttt{Model} sheet \emph{to the dollar}; the replica is the reference implementation
throughout.

% ================= 2 =================
\section{Model specification}

Let $O_t,H_t,L_t,C_t$ denote the daily open, high, low and close, and $F_t=H_t-L_t$ the range.

\subsection{Bayes sleeve: local linear trend}

The latent level $\ell_t$ and slope $b_t$ follow
\begin{equation}
\ell_t = \ell_{t-1} + b_{t-1} + \eta_t, \qquad
b_t = b_{t-1} + \zeta_t, \qquad
C_t = \ell_t + \varepsilon_t,
\end{equation}
with process and observation variances scaled by the realised range,
\begin{equation}
\operatorname{Var}(\eta_t)=(\phi_L F_t)^2,\qquad
\operatorname{Var}(\zeta_t)=(\psi F_t)^2,\qquad
\operatorname{Var}(\varepsilon_t)=(\lambda F_t)^2 .
\end{equation}
The Kalman recursion yields the one-step prediction and its uncertainty
\begin{equation}
\hat{\mu}_t = \ell_{t-1}+b_{t-1}, \qquad
\sigma_t = \sqrt{P^{11}_t+2P^{12}_t+P^{22}_t+(\phi_L F_t)^2+(\lambda F_t)^2}.
\end{equation}

\subsection{OU sleeve: windowed mean reversion}

Over a trailing window $W$ the sleeve estimates the mean $\bar{C}$, the AR(1) coefficient
$\rho$ (clipped to $[0,0.99]$) and the dispersion $s$, giving the forecast
\begin{equation}
\hat{\mu}^{\,\mathrm{OU}}_t = \bar{C}_t + \rho_t\left(C_{t-1}-\bar{C}_t\right).
\end{equation}

\subsection{Orders}

Each sleeve submits a bid capped by the open and by a fraction of the running peak $G_{t-1}$,
\begin{equation}
\beta_t=\min\!\Big(\hat{\mu}_t-k\,\sigma_t,\;\; O_t,\;\; G_{t-1}(1-c)\Big),
\label{eq:bid}
\end{equation}
and, once filled, rests a sell at the take-profit
\begin{equation}
\tau_t = \beta_t + \pi\,C_{t-1},
\label{eq:target}
\end{equation}
where $k$ is the bid buffer in standard deviations, $c$ the peak cap and $\pi$ the premium. A
position unrealised after $S$ calendar days is closed at the open (the stop).

% ================= 3 =================
\section{The fill-ordering problem}

A daily bar records $\{O_t,H_t,L_t,C_t\}$ but \emph{not the order in which $L_t$ and $H_t$
occurred}. The backtest books a same-day round trip whenever
\begin{equation}
L_t \le \beta_t \quad\text{and}\quad H_t \ge \tau_t ,
\label{eq:sameday}
\end{equation}
which presumes the low preceded the high. If the high in fact came first, the position could
not have been bought and sold that session: the trade is fictitious.

Condition~\eqref{eq:sameday} is not a minor edge case. Across the book it accounts for
\textbf{74\%} of Bayes trades and \textbf{90\%} of OU trades. Because returns compound, an
error concentrated in the fastest-recycling trades propagates multiplicatively.

Two sub-cases are decidable from daily data alone. A fill \emph{at the open}
($\beta_t \ge O_t$, i.e.\ the open was the binding cap in~\eqref{eq:bid}) is unambiguous: the
open is the session's first price, so any subsequent high necessarily follows the fill. A fill
on an intraday dip ($\beta_t<O_t$) is undecidable. Only $\sim\!15\%$ of same-day trades were
at-open fills, so daily bars alone leave the question overwhelmingly open.

% ================= 4 =================
\section{Verification methodology}

Let $\{(h_j,l_j)\}_{j=1}^{m}$ be the intraday bars of session $t$ in chronological order.
Define the first filling bar
\begin{equation}
j^{\ast} = \min\{\, j : l_j \le \beta_t \,\},
\end{equation}
and accept the same-day exit if and only if the target is reached at or after that bar,
\begin{equation}
\textsc{Real}(t) \iff \exists\, j \ge j^{\ast} \ \text{such that}\ h_j \ge \tau_t .
\label{eq:verify}
\end{equation}
Operationally this is evaluated with a suffix maximum $M_j=\max_{i\ge j} h_i$, so
$\textsc{Real}(t)\iff M_{j^\ast}\ge\tau_t$, at $O(m)$ preprocessing per session.

\subsection{Resolution adequacy}

Within a single bar the ordering of $l_j$ and $h_j$ remains unknown, so
predicate~\eqref{eq:verify} is optimistic at coarse resolution. Downsampling the NVDA
one-minute series establishes the tolerance:

\begin{table}[H]\centering\small
\renewcommand{\arraystretch}{1.15}
\begin{tabular}{lrrr}
\toprule
\textbf{Bar resolution} & \textbf{Annualised} & \textbf{Trades} & \textbf{Sharpe}\\
\midrule
1-minute (reference) & 36\% & 95 & 1.03\\
\rowcolor{bckeep} 5-minute & 36\% & 95 & 1.03\\
15-minute & 53\% & 111 & 1.55\\
\bottomrule
\end{tabular}
\caption*{\footnotesize Five-minute bars reproduce the one-minute result \emph{exactly}: with a
premium of $\pi\!\approx\!1.7\%$ the target is essentially never reached within five minutes of
the fill. Fifteen-minute bars inflate the result materially. Five-minute data is therefore
sufficient and was used for the remaining nine names.}
\end{table}

\subsection{Coverage}

Regular-hours bars only ($09{:}30$--$16{:}00$~ET); pre- and post-market bars are excluded as
the strategy trades the session. Coverage exceeds $94\%$ of sample sessions for every name.
Uncovered sessions fall back to the at-open rule, which is conservative.

% ================= 5 =================
\section{Results: fill rates and restated returns}

\begin{table}[H]\centering\small
\renewcommand{\arraystretch}{1.2}
\begin{tabular}{l r r r r}
\toprule
\textbf{Name} & \textbf{Fill rate} & \textbf{Backtest} & \textbf{Verified} & \textbf{Retained}\\
\midrule
\rowcolor{bckeep} RKLB & 69\% & 587\% & \textbf{113\%} & 0.19\\
PLTR & 64\% & 174\% & 43\% & 0.25\\
SPOT & 62\% & 149\% & 32\% & 0.21\\
TSM  & 59\% & 133\% & 53\% & 0.40\\
VRT  & 58\% & 219\% & 51\% & 0.23\\
VST  & 56\% & 257\% & 55\% & 0.21\\
TSLA & 54\% & 157\% & 18\% & 0.11\\
AVGO & 52\% & 143\% & 31\% & 0.22\\
NVDA & 50\% & 150\% & 36\% & 0.24\\
\rowcolor{bcbad} SOFI & 46\% & 231\% & 19\% & 0.08\\
\midrule
\textbf{Mean} & \textbf{57\%} & \textbf{$\sim$220\%} & \textbf{$\sim$45\%} & \textbf{0.21}\\
\bottomrule
\end{tabular}
\caption*{\footnotesize Fill rate is the proportion of same-day round trips satisfying
\eqref{eq:verify}. ``Retained'' is verified return divided by backtested return. Trade counts
fall by approximately $40\%$ on the verified basis.}
\end{table}

Three observations follow. First, the fill rate is \emph{tightly clustered} ($46\%$--$69\%$),
indicating a stable structural property rather than an idiosyncrasy of any one name. Second,
retention is uniformly poor ($0.08$--$0.40$): because compounding is multiplicative, losing
$43\%$ of the fastest-recycling trades costs far more than $43\%$ of the return. Third, the
ordering of names is broadly preserved --- RKLB remains the strongest by a wide margin --- so
relative judgements made on the old basis largely survive even though absolute levels do not.

\textbf{The strategy is sound; the backtest was measuring a version of it that fills better than
the market does.} A verified $\sim\!45\%$ annualised remains a strong result.

% ================= 6 =================
\section{Capital allocation between sleeves}

The sleeves are weakly correlated (mean daily correlation $\approx0.24$), so the split is a
genuine allocation decision. Earlier analysis on the \emph{unverified} basis favoured an
OU-heavy book. That conclusion inverts under verification, for a mechanical reason: the OU
sleeve depended on same-day round trips for $90\%$ of its trades against $74\%$ for Bayes, and
is therefore penalised disproportionately when fictitious fills are removed.

\begin{table}[H]\centering\small
\renewcommand{\arraystretch}{1.15}
\begin{tabular}{c r r r r}
\toprule
\textbf{Bayes share} & \textbf{Verified return} & \textbf{Sharpe} & \textbf{Max drawdown} & \textbf{Trades/yr}\\
\midrule
0\%   & 30\% & 0.64 & 39\% & unchanged\\
\rowcolor{bckeep} 50\% (current) & 45\% & 0.92 & 40\% & unchanged\\
70\%  & 50\% & 0.97 & 41\% & unchanged\\
\rowcolor{bcact} 80\%  & 53\% & 0.99 & 42\% & unchanged\\
100\% & 58\% & 1.02 & 43\% & reduced\\
\bottomrule
\end{tabular}
\caption*{\footnotesize Return and Sharpe both increase monotonically in the Bayes share for
only $\sim\!3$pp of additional drawdown. Trade count is invariant across interior splits --- the
split scales trade \emph{size}, not \emph{frequency} --- but falls at the corners, where one
sleeve is unfunded. A $5\%$ floor on the minor sleeve preserves the full count.
\textbf{These are in-sample figures and are ceilings, not forecasts}: the realised
out-of-sample gain is $+1.9$pp (\S6.1), so the planning figure is $\sim\!47\%$, not the $53\%$
or $58\%$ implied by the rows below. A per-name fitted split reaches $57\%$ in-sample but does
not survive validation.}
\end{table}

\subsection{Out-of-sample validation}

Three expanding folds per name, with the share chosen on the training window only and scored on
the unseen slice against the incumbent $50/50$:

\begin{table}[H]\centering\small
\renewcommand{\arraystretch}{1.15}
\begin{tabular}{lc}
\toprule
Folds won by the tilted allocation & \textbf{20 / 30}\\
Mean out-of-sample advantage & \textbf{$+1.9$pp}\\
Share selected on training data & $0.95$ in \textbf{30 of 30} folds\\
\bottomrule
\end{tabular}
\caption*{\footnotesize Perfect selection stability --- every training window on every name chose
the same share --- against roughly $p\approx0.05$ significance on the win rate. The in-sample
advantage of $+12$pp compresses to $+1.9$pp out of sample.}
\end{table}

Two qualifications matter. The validation supports a \emph{uniform} tilt, not the per-name
differentiation suggested by the full-sample fit: training windows selected a Bayes-heavy share
even for RKLB and TSLA, whose apparent OU preference is visible only with late data in view.
And the realised advantage is modest. A share of $70\text{--}80\%$ is therefore recommended
rather than the fitted $95\%$: the curve is flat across that range, and retaining a material OU
sleeve preserves decorrelation against regimes absent from this sample.

% ================= 7 =================
\section{Capital allocation across the book}

Verified returns differ substantially between names ($18\%$ to $113\%$), which invites weighting
by observed performance. That is sound only if performance persists. It does not.

\begin{table}[H]\centering\small
\renewcommand{\arraystretch}{1.15}
\begin{tabular}{lc}
\toprule
\textbf{Rank correlation between adjacent windows} & \textbf{$\rho$}\\
\midrule
Return, window $1\rightarrow2$ & $+0.71$\\
Return, window $2\rightarrow3$ & $-0.09$\\
Return, window $3\rightarrow4$ & $-0.30$\\
\midrule
\rowcolor{bckeep} Volatility, first half $\rightarrow$ second half & $\mathbf{+0.45}$\\
\bottomrule
\end{tabular}
\caption*{\footnotesize Return leadership does not persist and latterly reverses; volatility
does persist. This is the classical asymmetry --- risk is forecastable, return is not --- and it
determines which quantities may legitimately drive weights.}
\end{table}

Testing allocation schemes out of sample (weights formed on a training window, applied to the
next unseen window) confirms the implication:

\begin{table}[H]\centering\small
\renewcommand{\arraystretch}{1.15}
\begin{tabular}{l r r r r}
\toprule
\textbf{Scheme} & \textbf{Fold 1} & \textbf{Fold 2} & \textbf{Fold 3} & \textbf{Mean}\\
\midrule
\rowcolor{bckeep} Equal weight & 62.2\% & 8.6\% & 12.1\% & \textbf{27.6\%}\\
Inverse volatility & 53.6\% & 4.1\% & 15.3\% & 24.4\%\\
Inverse drawdown & 53.2\% & 3.2\% & 13.9\% & 23.4\%\\
Return weighted & 47.0\% & $-0.6\%$ & 7.2\% & 17.9\%\\
\rowcolor{bcbad} Risk-adjusted (return\,$\div$\,volatility) & 45.1\% & $-1.6\%$ & 9.9\% & \textbf{17.8\%}\\
\bottomrule
\end{tabular}
\caption*{\footnotesize Equal weighting dominates by roughly $10$pp annually. The
return-over-volatility rule currently implemented in the allocation sheet ranks \emph{last},
because it leans on the component that does not persist while correcting by the one that does.}
\end{table}

\textbf{Recommendation: weight the book equally}, or by inverse volatility if a risk tilt is
required (second best, and it uses only the persistent quantity). Do not weight by return or by
return over volatility. This is the largest single improvement identified in this work, and it
requires no change to the model.

% ================= 8 =================
\section{Register of rejected hypotheses}

Each hypothesis below produced a favourable in-sample result and was then tested on data unseen
during fitting. Six of seven did not survive.

\begin{table}[H]\centering\small
\renewcommand{\arraystretch}{1.25}
\begin{tabular}{p{5.1cm} p{4.4cm} c}
\toprule
\textbf{Hypothesis} & \textbf{Out-of-sample outcome} & \textbf{Verdict}\\
\midrule
Pre-market VWAP or actual open as the entry signal &
$-8$ to $-15$pp average; beats baseline for $0$--$2$ of $10$ names & \textcolor{bcbad}{\textbf{\color{red}Rejected}}\\
Re-optimising parameters for candidate names &
Frozen parameters win; refits gain $0$--$9$pp in-sample only & {\color{red}\textbf{Rejected}}\\
Wide premia fitted with same-day exits forbidden &
Premia $\sim\!5\times$ wider in-sample; underperforms on verified fills & {\color{red}\textbf{Rejected}}\\
Re-optimisation against verified fills (Bayes floor) &
$1/3$ folds; fitted values unstable across folds & {\color{red}\textbf{Rejected}}\\
Close-conditional premium (widen after a strong close) &
$14/24$ folds ($58\%$), $+2.7$pp; parameter pinned to the grid edge & {\color{red}\textbf{Rejected}}\\
Weekly cycle with forced week-end exit &
$6/18$ folds, mean $-10$pp against the frozen daily model & {\color{red}\textbf{Rejected}}\\
Per-name calendar stop in place of a uniform $50$ days &
$13/30$ folds, mean $-1.8$pp; selected stop scatters from $20$ to $150$ days & {\color{red}\textbf{Rejected}}\\
\rowcolor{bckeep} Tilt of capital toward the Bayes sleeve &
$20/30$ folds, $+1.9$pp, selection stable in $30/30$ & \textbf{Adopted}\\
\rowcolor{bckeep} Equal weighting across names &
Best of five schemes out of sample, $+9.8$pp over the incumbent rule & \textbf{Adopted}\\
\bottomrule
\end{tabular}
\end{table}

The stop-period result warrants a note, as an earlier study reported it more favourably. That
study fitted stops on the full sample and scored them on the unverified basis, obtaining $6$ of
$10$ names --- suggestive, and reported at the time as ``direction robust, magnitude not''. Under
the stricter protocol here (verified fills; stop chosen on training data only) it fails
outright, and the selected values contradict the earlier recommendation: that study proposed
$\sim\!40$ days generally with $60$ for RKLB, whereas training windows here select $90$--$150$
days for NVDA, TSM and VRT and $30$ days for RKLB. \textbf{The uniform $50$-day stop should be
retained}, and the earlier per-name table is superseded.

\subsection{Why the failures share a shape}

Every rejected hypothesis was a \emph{search}: an optimiser was asked to locate better values
within the existing machinery, and reliably located noise. Two diagnostics anticipated failure
in each case, and are recommended as standing tests:

\begin{enumerate}[leftmargin=1.5em,itemsep=2pt]
  \item \kpi{Parameter instability.}{} Fitted values that swing widely between folds --- the
        weekly model's premium ranged $0.06$ to $0.20$ on a single name --- indicate the
        optimiser is fitting each window's noise. The adopted allocation, by contrast, was
        selected identically in all thirty folds.
  \item \kpi{Boundary seeking.}{} An optimum that migrates to the edge of its search range and
        does not turn over is not an optimum. This occurred in the wide-premium, close-conditional
        and weekly experiments.
\end{enumerate}

The two adopted changes share the opposite character. Neither fits a value: each follows from a
\emph{measured, persistent} property of the data --- that OU depends more heavily than Bayes on
unverifiable fills, and that return leadership does not persist while volatility does. Both were
selected identically across folds. The distinction between fitting a parameter and exploiting a
measured regularity is, on this evidence, what separates changes that survive from those that
do not.

% ================= 9 =================
\section{Conclusions and recommendations}

\begin{enumerate}[leftmargin=1.5em,itemsep=3pt]
  \item \kpi{Weight the book equally.}{} Replace the return-over-volatility rule with equal
        weights (or inverse volatility if a risk tilt is wanted). Worth $\sim\!+10$pp annually
        out of sample --- the largest improvement available, and it requires no model change.
  \item \kpi{Restate expected performance.}{} Plan against $\sim\!45\%$ annualised at the
        current sleeve split, rising to $\sim\!47\%$ once the validated tilt is applied, with
        approximately $40\%$ fewer trades than the daily backtest reports. The prior figure of
        $\sim\!220\%$ is not attainable, being contingent on same-day fills that occur roughly
        $57\%$ of the time. In-sample split optima of $53\text{--}58\%$ are ceilings, not
        forecasts.
  \item \kpi{Adopt a moderate Bayes tilt.}{} Move the sleeve allocation to $70\text{--}80\%$
        Bayes, uniformly, with a $5\%$ floor on the minor sleeve. Expected benefit
        $\sim\!+2$pp annualised at no cost to trade frequency.
  \item \kpi{Leave all model parameters unchanged, including the uniform $50$-day stop.}{}
        Seven independent attempts to improve them failed out of sample. The edge is structural,
        not parametric.
  \item \kpi{Adopt the verified basis as standard.}{} Any future candidate name or model change
        should be evaluated on intraday-verified fills. Five-minute resolution is sufficient;
        fifteen-minute is not.
  \item \kpi{Re-rank prior candidate studies.}{} Assessments of new names (CCL, MU, CVNA, LLY
        and the diversifier set) were conducted on the unverified basis. Since fill rates vary
        between $46\%$ and $69\%$, their relative ordering may shift once restated.
\end{enumerate}

\textbf{Principal caveat.} The sample spans April 2024 to mid-2026, a sustained bull regime for
high-beta US equities. No conclusion here is tested against a sideways or declining market.
This is the primary reason to retain a material OU sleeve despite its weaker verified
performance: its value is decorrelation in regimes the data does not contain.

\vfill
{\footnotesize\color{bcgrey} Data: daily OHLC April 2024--July 2026 for ten names; one-minute
bars for NVDA and five-minute bars for the remaining nine, regular hours. The Python replica
reproduces every cached workbook sheet to the dollar. All out-of-sample results use expanding
walk-forward with parameters selected on training windows only.}

\end{document}
